\section{Round-seventeen positive closure: one global history kernel and forced compressed memory}
\label{sec:r17-a4}

The history dynamics is constructed globally from the legal return graph.  No
partition of unity is used to mix local kernels.  Weak-Harris contraction is
proved in a patched distance from a derived drift, a derived small-set
minorization, and the summable-variation coupling.  The memory identity is a
common-domain block-resolvent theorem and retains the unresolved initial
forcing.

\subsection{The global past kernel}

Let \(\Omega^-\) be the space of legal left-infinite return-edge histories
with the current suspension residual.  Put
\[
 d_\beta(x,y)=\sum_{j\ge0}e^{-\beta j}
   \bigl(1\wedge d_E(x_{-j},y_{-j})\bigr),
 \qquad
 W(x)=\exp\{\eta(r(x_0)+\tau(x_0)+\mathfrak s(x_0))\}.
\]
The reference Gibbs conditional probability of the next edge is denoted
\(g(e\mid x)\).  Summable variations give
\[
 \sum_{n\ge1}\operatorname{var}_n(\log g)<\infty.
\]
Appending a legal edge and shifting the past defines one global Markov kernel
\(P\) on \(\Omega^-\).

\begin{lemma}[Derived one-step coupling]
\label{lem:r17-a4-coupling}
There are \(C,c>0\) such that, if two histories agree through depth \(N\),
then their next-edge kernels have total variation distance at most
\(Ce^{-cN}\).  Under a maximal coupling, a successful match multiplies
\(d_\beta\) by \(e^{-\beta}\); a failed match has probability at most
\(Cd_\beta(x,y)^\alpha\) for some \(\alpha>0\).
\end{lemma}

\begin{proof}
The Gibbs ratio of two next-edge probabilities is the exponential of the
variation of the induced potential on histories sharing their last \(N\)
edges.  Summable variation bounds the logarithmic ratio by \(Ce^{-cN}\), and
Pinsker's elementary two-measure estimate gives the total variation bound.
Choose \(N\) comparable with \(|\log d_\beta|\).  If the next edges match, all
old discrepancies shift one coordinate into the past, giving the factor
\(e^{-\beta}\).
\end{proof}

\begin{lemma}[Lyapunov drift and a common small set]
\label{lem:r17-a4-drift}
For sufficiently small \(\eta>0\),
\[
 PW\le\lambda W+C,
 \qquad \lambda<1.
\]
For every \(R<\infty\), the set \(\{W\le R\}\) is \(n_R\)-small: there are a
probability \(\nu_R\) and \(\epsilon_R>0\) such that
\[
 P^{n_R}(x,\cdot)\ge\epsilon_R\nu_R(\cdot)
 \quad\text{for all }W(x)\le R.
\]
\end{lemma}

\begin{proof}
The exponential return/roof/singularity moment of the induced Gibbs kernel
implies a Foster--Lyapunov estimate after reducing \(\eta\); the current edge
is shifted out and the new exponential moment is uniformly bounded.  On a
Lyapunov sublevel only finitely many edge cylinders carry all but a fixed
small probability.  The finite connector property supplies a common legal
word from each such cylinder to one fixed magnet cylinder.  Bounded Gibbs
distortion gives a uniform positive lower bound on the probability of that
word and on a common terminal subcylinder, proving minorization.
\end{proof}

Define \(d_\epsilon=1\wedge\epsilon^{-1}d_\beta^\alpha\) and
\[
 \widetilde d(x,y)=
 \sqrt{d_\epsilon(x,y)
       \bigl(1+\delta W(x)+\delta W(y)\bigr)}.
\]

\begin{theorem}[Global weak-Harris spectral theorem]
\label{thm:r17-a4-harris}
For suitable \(\epsilon,\delta\) there are \(n_0\) and \(\rho<1\) such that
\[
 \mathcal W_{\widetilde d}(P^{n_0}(x,\cdot),P^{n_0}(y,\cdot))
 \le\rho\widetilde d(x,y).
\]
Consequently \(P\) has a unique invariant probability \(\nu\), a spectral gap
on the weighted Lipschitz space
\[
 \mathcal B=\{f:\|f\|_W+\operatorname{Lip}_{\widetilde d}f<\infty\},
\]
and exponentially summable correlations.  This kernel is the original global
Doob/history kernel; no chart mixture has been introduced.
\end{theorem}

\begin{proof}
For small \(d_\epsilon\), use the coupling of
Lemma~\ref{lem:r17-a4-coupling}; the failed-match term is absorbed by choosing
\(\epsilon\) small.  For large distance but bounded \(W\), use the common
small-set coupling from Lemma~\ref{lem:r17-a4-drift}.  Outside that set, the
Lyapunov drift contracts the weight factor.  The three estimates are exactly
the weak-Harris contraction in the patched distance.  Iteration gives the
spectral gap and correlation bound on \(\mathcal B\).
\end{proof}

\subsection{One fixed multiplier algebra and Doob tower}

Fix \(0<\eta_0<\eta\) and define
\[
 \mathcal A=\left\{V:
  \|V\|_{\mathcal A}:=
  \sup_x{|V(x)|\over1+\log W(x)}
  +\operatorname{Lip}_{\widetilde d}V<\infty\right\}.
\]
Let \(\mathcal B_\theta\) be the same weighted Lipschitz space with weight
\(W^\theta\), \(0\le\theta\le1\).

\begin{theorem}[Fixed analytic Feynman--Kac chart]
\label{thm:r17-a4-fk}
There is \(r_0>0\) and fixed \(0<\theta<\theta'<1\) such that
\[
 P_Vf=P(e^Vf)
\]
is an analytic family in \(\|V\|_{\mathcal A}<r_0\) from
\(\mathcal B_\theta\) to itself and from
\(\mathcal B_\theta\) to \(\mathcal B_{\theta'}\).  It has a simple positive
eigenvalue \(e^{\Lambda(V)}\), a positive eigenfunction \(h_V\), and a dual
eigenmeasure \(\nu_V\).  The normalized kernel
\[
 \widehat P_Vf=e^{-\Lambda(V)}h_V^{-1}P_V(h_Vf)
\]
is Markov and satisfies the exact nonlinear tower
\[
 \mathcal E_{m+n}^V F
 =\mathcal E_m^V(\mathcal E_n^VF).
\]
\end{theorem}

\begin{proof}
The inequality \(|V|\le r_0(1+\log W)\) implies
\(e^{|V|}W^\theta\le C W^{\theta'}\) when \(r_0<\theta'-\theta\).
The Lipschitz product estimate and
Theorem~\ref{thm:r17-a4-harris} give a uniform bounded multiplier map on the
fixed pair of spaces.  The exponential power series converges in operator
norm, hence is analytic.  Analytic perturbation of the isolated positive
eigenvalue gives \(\Lambda,h_V,\nu_V\).  The displayed Doob transform is
Markov by the eigenfunction equation, and its ordinary Markov tower becomes
the nonlinear tower after logarithmic conjugation.
\end{proof}

\subsection{Martingale and rough-path limit}

For centered \(g\in\mathcal B\), the Poisson series
\(h=\sum_{n\ge0}P^ng\) converges in \(\mathcal B\), and
\[
 d_{n+1}=h(X_{n+1})-Ph(X_n)
\]
is a martingale difference.

\begin{theorem}[Enhanced invariance principle]
\label{thm:r17-a4-rough}
For every finite vector of centered observables in \(\mathcal B\) with a
slightly higher \(W\)-moment, the suspension additive functional and its
iterated integral converge jointly to a Brownian rough path.  The covariance
is the absolutely convergent Green--Kubo series and the antisymmetric area
drift is the convergent martingale--coboundary area series.  The conclusion is
stable under conditioning on a compact past cylinder.
\end{theorem}

\begin{proof}
The spectral gap gives the Poisson solution and an \(L^{2+\epsilon}\)
martingale--coboundary decomposition.  The martingale functional CLT and its
iterated-integral version apply; the coboundary contributes only the explicit
area correction.  Exponential roof moments control suspension inversion and
the final residual roof.  Conditioning on a compact past cylinder changes the
initial law by a uniformly bounded density, so the same Lindeberg and bracket
estimates hold.
\end{proof}

\subsection{Raw renewal resolvent}

Let \(R_V(z)\) be the first-return operator weighted by
\(e^{-z\tau+V}\).  Its matrix coefficients are the Laplace--Fourier transforms
controlled by A2.

\begin{theorem}[Unsmoothed suspension resolvent]
\label{thm:r17-a4-renewal}
For \(V\) in a compact subball of \(\mathcal A\), the renewal identity
\[
 (z-L_V)^{-1}=A_V(z)+B_V(z)(I-R_V(z))^{-1}C_V(z)
\]
holds from a common graph domain \(D\) to \(L^2(\nu_V)\).  The terms are
meromorphic in \(\Re z>-\gamma\), and after the finitely many resonances are
removed their matrix coefficients and first two source derivatives are
integrable on vertical lines.  The bounds are uniform in \(V\).
\end{theorem}

\begin{proof}
Decompose a suspension path into its initial residual, complete returns, and
final residual.  Fubini on the common exponential roof weight gives the
renewal identity first for \(\Re z\) large.  The A2 trace-free matrix-
coefficient Fourier estimate gives an integrable bound for the complete-return
series; the residual operators are bounded by the exponential roof moment.
Analytic Fredholm reduction on the isolated finite-dimensional leading
spectral subspace gives the meromorphic continuation.  The same estimates
hold after two source derivatives because the multiplier algebra is fixed.
\end{proof}

\subsection{Common-domain forced compressed memory}

Let \(H_V=L^2(\nu_V)\), let \(L_V\) be the generator of the normalized
suspension semigroup, and let \(P_V\) be the orthogonal projection onto a
finite-dimensional resolved space \(\mathcal R_V\subset D(L_V^2)\).  The
spaces are transported to one reference Hilbert space by the square-root
density isometry.  Put \(Q_V=I-P_V\) and
\[
 C_V(z)=P_V(z-L_V)^{-1}P_V.
\]

\begin{lemma}[Right-half-plane compression bound]
\label{lem:r17-a4-compression}
For \(\Re z\ge\sigma>0\), \(C_V(z)\) is invertible on \(\mathcal R_V\).  On
compact vertical contour segments its least singular value is uniformly
positive, while for large \(|z|\),
\[
 C_V(z)=z^{-1}P_V+O(|z|^{-2}).
\]
\end{lemma}

\begin{proof}
The normalized Markov semigroup is an \(L^2(\nu_V)\)-contraction, so \(L_V\)
is dissipative.  If \(x\in\mathcal R_V\) and
\(y=(z-L_V)^{-1}x\), then
\[
 \Re\langle C_V(z)x,x\rangle
 =\Re\langle y,(z-L_V)y\rangle
 \ge\sigma\|y\|^2.
\]
Thus \(C_V(z)x=0\) implies \(x=0\); finite dimensionality gives
invertibility.  Compactness of the source set and contour gives the uniform
minimum singular value.  The large-\(z\) expansion follows from
\(\mathcal R_V\subset D(L_V^2)\).
\end{proof}

Define
\[
 \widehat K_V(z)=zP_V-P_VL_VP_V-C_V(z)^{-1}.
\]

\begin{theorem}[Forced Mori--Zwanzig identity]
\label{thm:r17-a4-memory}
For \(A\in D(L_V)\), the resolved trajectory
\(x_A(t)=P_Ve^{tL_V}A\) satisfies, in \(\mathcal R_V\),
\[
 \dot x_A(t)=P_VL_VP_Vx_A(t)
 +\int_0^tK_V(t-s)x_A(s)\,ds+\eta_A(t),
\]
where
\[
 \widehat\eta_A(z)=C_V(z)^{-1}
 P_V(z-L_V)^{-1}Q_VA.
\]
The kernel is the causal inverse Laplace transform of
\(\widehat K_V\).  After adjoining the finitely many left-half-plane
transmission-zero residue modes, the residual kernel and forcing are
exponentially integrable.  If \(A\in\mathcal R_V\), the forcing vanishes.
\end{theorem}

\begin{proof}
The Laplace transform of \(x_A\) is
\[
 X_A(z)=C_V(z)P_VA+P_V(z-L_V)^{-1}Q_VA.
\]
Multiply by \(C_V(z)^{-1}\) and rearrange to obtain the transform of the
stated Volterra equation, including its initial value and forcing.  All
products are defined because \(\mathcal R_V\subset D(L_V^2)\).
Lemma~\ref{lem:r17-a4-compression} gives the inverse on the right half-plane.
Theorem~\ref{thm:r17-a4-renewal} gives meromorphic continuation.  A zero of the
continued compression is promoted, with its full Riesz residue range, to the
resolved space; the Schur complement then has no such pole.  Finitely many
steps remove all zeros in a fixed strip, and contour shifting gives
exponential integrability.
\end{proof}

\begin{theorem}[Round-seventeen A4 closure]
\label{thm:r17-a4-main}
The prepared Sinai history has one global weak-Harris kernel, one fixed
analytic Feynman--Kac algebra, a conditional enhanced invariance principle, an
unsmoothed renewal resolvent, and a common-domain forced compressed-memory
equation with quantitative inverse bounds.
\end{theorem}

\begin{proof}
Combine Theorems~\ref{thm:r17-a4-harris},
\ref{thm:r17-a4-fk}, \ref{thm:r17-a4-rough},
\ref{thm:r17-a4-renewal}, and \ref{thm:r17-a4-memory}.
\end{proof}
